\chapter{Probabilité à grande vitesse}\label{ch:qm:probability-at-speed}

À quatre heures moins dix, le déséquilibre de l'enchère de clôture apparaît à l'écran, et l'estimation courante que le desk
fait du cours de clôture, un chiffre qu'il révise depuis le matin, bondit de 14
centimes. Un chercheur pose une question plus fine que celle de savoir si l'estimation était juste~: ses révisions auraient-elles pu
être prédites à partir des révisions qui les ont précédées~? Si oui, l'estimation n'utilisait pas ce qu'elle savait, et quiconque aurait
opéré contre elle aurait gagné de l'argent. Une prévision honnête
d'une quantité fixe est une \omterm{def:qm:probability-at-speed:condexp}{espérance conditionnelle} le long d'un flux d'information, et un tel processus possède
une propriété caractéristique~: ses révisions ne peuvent pas être prévues. Ce chapitre présente
au rythme d'un lecteur qui a déjà rencontré la théorie de la mesure, les objets que calcule la série
(l'information, l'\omterm{def:qm:probability-at-speed:condexp}{espérance conditionnelle}, les \omterm{def:qm:probability-at-speed:martingale}{martingales}, les \omterm{def:qm:probability-at-speed:stopping}{temps d'arrêt}, les \omterm{def:qm:probability-at-speed:change}{changements de probabilité})
et se conclut en fixant les notations qu'emploie chaque livre suivant.

\section{Information~: tribus, filtrations et espérance conditionnelle}

Un espace de probabilité $(\Omega, \mathcal F, \P)$ est supposé connu. Ce que la finance ajoute,
c'est le temps~: à chaque instant, certains événements sont tranchés, et les événements tranchés s'accumulent.

\begin{definition}[Filtration, processus adapté]\label{def:qm:probability-at-speed:filtration}
Une \emph{filtration}\index{filtration} est une famille $\mathbb F = (\mathcal F_t)_{t \in
\mathbb T}$ de sous-$\sigma$-algèbres de $\mathcal F$, croissante en $t$~: $\mathcal F_s
\subseteq \mathcal F_t$ pour $s \le t$. L'ensemble d'indices $\mathbb T$ vaut $\{0, 1, \dots, n\}$ ou
$[0, T]$~; en temps continu, la filtration est supposée continue à droite et complète (les
\emph{conditions habituelles}). Un processus $(X_t)$ est un \emph{processus adapté}\index{processus adapté}
si chaque $X_t$ est $\mathcal F_t$-mesurable~: sa valeur en $t$ est connue en $t$.
\end{definition}

$\mathcal F_t$ est ce qui est connu en $t$~: les cours publiés jusqu'ici, les cartes retournées. Une
règle de trading doit être adaptée~; un backtest qui ne l'est pas a regardé dans le futur (chapitre~3).

\begin{definition}[Espérance conditionnelle]\label{def:qm:probability-at-speed:condexp}
Soient $X$ intégrable et $\mathcal G \subseteq \mathcal F$ une sous-$\sigma$-algèbre. On définit
l'\emph{espérance conditionnelle}\index{espérance conditionnelle} $\E[X \mid \mathcal G]$ comme
la variable aléatoire presque sûrement unique $\mathcal G$-mesurable et intégrable $Y$ telle que
$\E[X\mathbf 1_A] = \E[Y\mathbf 1_A]$ pour tout $A \in \mathcal G$. On écrit $\E_t[X] =
\E[X \mid \mathcal F_t]$.
\end{definition}

L'existence découle du théorème de Radon--Nikodym appliqué à la mesure $A \mapsto \E[X\mathbf
1_A]$ sur $\mathcal G$. Pour un $X$ de carré intégrable, on dispose d'une image plus utile~:
$\E[X \mid \mathcal G]$ est la projection orthogonale de $X$ sur le sous-espace fermé
$L^2(\mathcal G)$, la meilleure prévision de $X$ au sens de l'erreur quadratique moyenne parmi toutes les variables aléatoires $\mathcal
G$-mesurables. Tout ce qu'un prévisionniste peut calculer à partir de l'information
$\mathcal G$ se trouve dans ce sous-espace, et l'erreur de prévision $X - \E[X \mid \mathcal G]$ est
orthogonale à toutes ces variables.

\begin{proposition}[Règles de l'espérance conditionnelle]\label{prop:qm:probability-at-speed:rules}
Pour $X, Y$ intégrables et $\mathcal H \subseteq \mathcal G$~:
(i) \emph{tour}~: $\E[\E[X \mid \mathcal G] \mid \mathcal H] = \E[X \mid \mathcal H]$~;
(ii) \emph{sortie de ce qui est connu}~: si $Y$ est $\mathcal G$-mesurable et $XY$ intégrable,
$\E[XY \mid \mathcal G] = Y\E[X \mid \mathcal G]$~;
(iii) \emph{indépendance}~: si $X$ est indépendante de $\mathcal G$, $\E[X \mid \mathcal G] =
\E[X]$~;
(iv) \emph{Jensen}~: pour $\varphi$ convexe avec $\varphi(X)$ intégrable, $\varphi(\E[X \mid
\mathcal G]) \le \E[\varphi(X) \mid \mathcal G]$.
\end{proposition}

\begin{proof}
Vérifier les intégrales définissantes~: (i) pour $A \in \mathcal H \subseteq \mathcal G$, $\E[\E[X \mid
\mathcal G]\mathbf 1_A] = \E[X\mathbf 1_A]$~; (ii) pour $Y = \mathbf 1_B$, $B \in \mathcal G$,
puis passer à la limite~; (iii) $\E[X\mathbf 1_A] = \E[X]\P(A)$~; (iv) $\varphi$ est un supremum de
fonctions affines en nombre dénombrable.
\end{proof}

La \Cref{fig:qm:probability-at-speed:tree} montre les règles à l'œuvre sur le plus petit exemple
qui les réunit toutes. Un prix part de 100 et monte ou descend d'un dollar à chacun des deux
jours, avec probabilités égales~; $X = (S_2 - 100)^+$. Après un jour, l'information est
celle de la moitié de l'arbre dans laquelle on se trouve, et $\E_1[X]$ fait la moyenne de $X$ sur cette moitié.

\begin{omfigure}
\begin{tikzpicture}[x=3.2cm,y=0.95cm,font=\small,
  nd/.style={draw=omDef,rounded corners=2pt,fill=omDef!6,inner sep=3pt,align=center}]
\node[nd] (a) at (0,0) {$S_0 = 100$\\$\E_0[X] = 0.5$};
\node[nd] (b) at (1,1.6) {$S_1 = 101$\\$\E_1[X] = 1$};
\node[nd] (c) at (1,-1.6) {$S_1 = 99$\\$\E_1[X] = 0$};
\node[nd] (d) at (2,2.4) {HH: $S_2 = 102$, $X = 2$};
\node[nd] (e) at (2,0.8) {HT: $S_2 = 100$, $X = 0$};
\node[nd] (f) at (2,-0.8) {TH: $S_2 = 100$, $X = 0$};
\node[nd] (g) at (2,-2.4) {TT: $S_2 = 98$, $X = 0$};
\draw[-{Stealth},thick,gray] (a) -- (b); \draw[-{Stealth},thick,gray] (a) -- (c);
\draw[-{Stealth},thick,gray] (b) -- (d.west); \draw[-{Stealth},thick,gray] (b) -- (e.west);
\draw[-{Stealth},thick,gray] (c) -- (f.west); \draw[-{Stealth},thick,gray] (c) -- (g.west);
\draw[omThm,dashed,rounded corners] (2.72,3.0) rectangle (3.02,0.2);
\draw[omThm,dashed,rounded corners] (2.72,-0.2) rectangle (3.02,-3.0);
\node[omThm,font=\footnotesize,align=left,anchor=west] at (3.05,1.6) {un atome\\de $\mathcal F_1$};
\node[omThm,font=\footnotesize,align=left,anchor=west] at (3.05,-1.6) {un atome\\de $\mathcal F_1$};
\node[font=\footnotesize,gray] at (0,-3.4) {$t = 0$};
\node[font=\footnotesize,gray] at (1,-3.4) {$t = 1$};
\node[font=\footnotesize,gray] at (2,-3.4) {$t = 2$};
\end{tikzpicture}

\omcaption{L'\omterm{def:qm:probability-at-speed:condexp}{espérance conditionnelle} comme moyenne sur ce qui reste inconnu. $\mathcal F_1$
découpe les quatre issues en deux atomes (en pointillés)~; $\E_1[X]$ est la moyenne de $X$ sur l'atome
atteint, et $\E_0[X]$ est la moyenne de $\E_1[X]$~: la règle de la tour. Les trois valeurs
$0.5 \to 1 \to 2$ le long du chemin supérieur forment une \omterm{def:qm:probability-at-speed:doob}{martingale de Doob}.}\label{fig:qm:probability-at-speed:tree}
\end{omfigure}

Les lois se traitent par leurs transformées. Celle que la série utilise le plus est définie ici, afin
que les livres suivants (processus de Lévy au chapitre~6, tarification par transformée de Fourier au chapitre~28, modèles de volatilité
stochastique dans le livre~5) puissent y renvoyer.

\begin{definition}[Fonction caractéristique]\label{def:qm:probability-at-speed:charfn}
La \emph{fonction caractéristique}\index{fonction caractéristique} d'une variable aléatoire $X$
est $\varphi_X(u) = \E[e^{\iu uX}]$, $u \in \R$~; pour un vecteur aléatoire, $\varphi_X(u) =
\E[e^{\iu u^\top X}]$, $u \in \R^d$. Elle existe pour toute loi, la détermine, et transforme les sommes
de variables indépendantes en produits.
\end{definition}

\begin{theorem}[Théorème central limite]\label{thm:qm:probability-at-speed:clt}
Soient $X_1, X_2, \dots$ des variables indépendantes d'espérances $\mu_i$ et de variances $\sigma_i^2$, $s_n^2 =
\sum_{i \le n}\sigma_i^2$, et supposons que la condition de Lindeberg $s_n^{-2}\sum_{i\le n}\E[(X_i -
\mu_i)^2\mathbf 1_{\{|X_i - \mu_i| > \varepsilon s_n\}}] \to 0$ soit vérifiée pour tout $\varepsilon >
0$ (elle l'est pour des variables identiquement distribuées de variance finie). Alors
$s_n^{-1}\sum_{i \le n}(X_i - \mu_i) \xrightarrow{d} \mathcal N(0, 1)$.
\end{theorem}

\admitted

La démonstration développe la \omterm{def:qm:probability-at-speed:charfn}{fonction caractéristique} à l'ordre deux et s'appuie sur le
théorème de continuité de Lévy (Williams, 1991). Toutes les erreurs types de la partie~II reposent sur lui.

\section{Martingales et martingale de Doob d'une prévision}

\begin{definition}[Martingale]\label{def:qm:probability-at-speed:martingale}
Un \omterm{def:qm:probability-at-speed:filtration}{processus adapté} et intégrable $(M_t)$ est une \emph{martingale}\index{martingale} si
$\E_s[M_t] = M_s$ pour tout $s \le t$~; une \emph{sous-martingale}\index{sous-martingale} si
$\E_s[M_t] \ge M_s$~; une \emph{surmartingale}\index{surmartingale} si $\E_s[M_t] \le
M_s$.
\end{definition}

Une \omterm{def:qm:probability-at-speed:martingale}{martingale} est un jeu équitable vu de n'importe quelle date. Quatre exemples reviennent dans ce livre. Avec des variables $\xi_i$ indépendantes, prenant les valeurs $\pm 1$ avec
probabilités $p$ et $q = 1 - p$, et $S_n = \sum_{i \le n}\xi_i$~:
\begin{itemize}
\item $S_n$ est une \omterm{def:qm:probability-at-speed:martingale}{martingale} lorsque $p = \tfrac12$, une \omterm{def:qm:probability-at-speed:martingale}{sous-martingale} lorsque $p > \tfrac12$~;
\item $S_n^2 - n$ est une \omterm{def:qm:probability-at-speed:martingale}{martingale} lorsque $p = \tfrac12$~: la variance croît d'une unité à chaque pas~;
\item $(q/p)^{S_n}$ est une \omterm{def:qm:probability-at-speed:martingale}{martingale} pour tout $p$, puisque $\E[(q/p)^{\xi}] = p\,q/p + q\,p/q = 1$~;
\item $\exp(\theta S_n - n\ln\cosh\theta)$ est une \omterm{def:qm:probability-at-speed:martingale}{martingale} lorsque $p = \tfrac12$, pour tout réel
$\theta$~: la \omterm{def:qm:probability-at-speed:martingale}{martingale} exponentielle, prototype des \omterm{def:qm:probability-at-speed:change}{processus densité} de
la \cref{sec:qm:probability-at-speed:measure}.
\end{itemize}

L'exemple le plus important pour un desk de trading n'est pas un prix mais une prévision.

\begin{definition}[Martingale de Doob]\label{def:qm:probability-at-speed:doob}
Pour une variable aléatoire intégrable $X$ et une \omterm{def:qm:probability-at-speed:filtration}{filtration} $\mathbb F$, la \emph{martingale
de Doob}\index{martingale de Doob} de $X$ est $M_t = \E_t[X]$.
\end{definition}

C'est une \omterm{def:qm:probability-at-speed:martingale}{martingale} par la loi des espérances itérées. Toute prévision honnête d'une quantité fixe (un cours
de clôture, la somme de cinq cartes, l'indice d'inflation publié le mois suivant) est une \omterm{def:qm:probability-at-speed:doob}{martingale de Doob} le long de la
\omterm{def:qm:probability-at-speed:filtration}{filtration} du prévisionniste, et ce qui la rend testable, c'est que ses révisions sont orthogonales.

\begin{proposition}[Accroissements orthogonaux]\label{prop:qm:probability-at-speed:orthogonal}
Soit $(M_k)_{k = 0,\dots,n}$ une \omterm{def:qm:probability-at-speed:martingale}{martingale} de carré intégrable d'accroissements $D_k = M_k -
M_{k-1}$. Alors $\E[D_k Y] = 0$ pour toute variable $\mathcal F_{k-1}$-mesurable et de carré
intégrable $Y$~; en particulier $\E[D_jD_k] = 0$ pour $j < k$, et
\[
\Var(M_n) = \Var(M_0) + \sum_{k=1}^n \E[D_k^2].
\]
\end{proposition}

\begin{proof}
$\E[D_kY] = \E[Y\E_{k-1}[D_k]] = 0$ par sortie de ce qui est connu. Avec $Y = D_j$, $j < k$, les
termes croisés de $(M_n - M_0)^2 = (\sum_k D_k)^2$ s'annulent, et $\Cov(M_0, D_k) = 0$ de même.
\end{proof}

Deux conséquences sous-tendent le tutoriel et le problème de ce chapitre. D'abord, une prévision dont
les révisions peuvent être prédites à partir du passé (une régression de $D_k$ sur $D_{k-1}$ avec une
pente non nulle) n'est pas une \omterm{def:qm:probability-at-speed:condexp}{espérance conditionnelle}~: elle laisse de l'information inutilisée. Ensuite, la
variance de la valeur finale se décompose en variances des révisions, si bien qu'on peut dire quelle
part de l'incertitude d'un cours de clôture est levée à chaque moment de la journée
(la \cref{fig:qm:probability-at-speed:day}).

\begin{omfigure}
\begin{tikzpicture}
\begin{groupplot}[group style={group size=2 by 1,horizontal sep=1.6cm},
  width=7.6cm,height=4.8cm,xlabel={minutes après 9~h~30},
  tick label style={font=\small},label style={font=\small},
  xmin=0,xmax=395,grid=major,grid style={gray!25},xtick={0,120,240,380},
  table/col sep=comma]
\nextgroupplot[ylabel={prévision du cours de clôture (USD)},ymin=99.2,ymax=102.4]
\addplot[omDef,thick] table[x=minute,y=forecast]{figdata/methods/01-probability-at-speed/day.csv};
\draw[omThm,dashed] (axis cs:380,99.2) -- (axis cs:380,102.4);
\node[font=\footnotesize,anchor=east,omThm] at (axis cs:374,99.5) {annonce de 15~h~50~: $+0.48$};
\nextgroupplot[ylabel={part de Var(clôture) résolue},ymin=0,ymax=1.05,
  legend style={font=\footnotesize,at={(0.03,0.97)},anchor=north west,fill=white,draw=none}]
\addplot[omDef,very thick] table[x=minute,y=theory]{figdata/methods/01-probability-at-speed/resolved.csv};
\addplot[omThm,only marks,mark=*,mark size=1.3pt] table[x=minute,y=sim]{figdata/methods/01-probability-at-speed/resolved.csv};
\legend{formule,20\,000 jours}
\end{groupplot}
\end{tikzpicture}

\omcaption{À gauche~: une journée simulée de $M_t = \E_t[\text{clôture}]$ dans le
modèle du problème du week-end (une marche aléatoire d'écart-type quotidien de 1.80~\$ sur une action à 100~\$, et des annonces
sur le déséquilibre de clôture, d'écart-type 0.30~\$, à 15~h~50, $+0.48$ ce jour-là).
À droite~: la part de la variance de la clôture résolue à chaque moment de la journée, selon
la \cref{prop:qm:probability-at-speed:orthogonal}~: linéaire sur toute la séance, un saut de 2.7~\% à l'annonce, les
2.5~\% restants dans les dix dernières minutes. Données~: le tutoriel du chapitre,
graine fixée.}\label{fig:qm:probability-at-speed:day}
\end{omfigure}

L'inégalité maximale borne jusqu'où une \omterm{def:qm:probability-at-speed:martingale}{martingale} peut s'écarter sur une période, à partir de sa valeur finale~;
le chapitre~2 l'utilise pour contrôler les trajectoires browniennes.

\begin{proposition}[Inégalité maximale de Doob]\label{prop:qm:probability-at-speed:doob}
Pour une \omterm{def:qm:probability-at-speed:martingale}{sous-martingale} à valeurs positives $(X_k)_{k \le n}$ et $c > 0$, $c\,\P(\max_{k\le n}X_k \ge c)
\le \E[X_n\mathbf 1_{\{\max_k X_k \ge c\}}] \le \E[X_n]$~; et pour $p > 1$, $\E[\max_{k \le
n}X_k^p] \le (p/(p-1))^p\E[X_n^p]$. Pour une \omterm{def:qm:probability-at-speed:martingale}{martingale} de carré intégrable, appliquée à $|M_k|$~:
$\E[\max_{k\le n}M_k^2] \le 4\E[M_n^2]$.
\end{proposition}

\admitted

\section{Temps d'arrêt et arrêt optionnel}

\begin{definition}[Temps d'arrêt]\label{def:qm:probability-at-speed:stopping}
Un temps aléatoire $\tau$ à valeurs dans $\mathbb T \cup \{\infty\}$ est un \emph{temps
d'arrêt}\index{temps d'arrêt} si $\{\tau \le t\} \in \mathcal F_t$ pour tout $t$~: le fait
qu'il se soit produit au plus tard en $t$ est connu en $t$.
\end{definition}

«~Le premier instant où le cours touche 101~» est un \omterm{def:qm:probability-at-speed:stopping}{temps d'arrêt}~; «~l'instant du plus haut de la journée~»
n'en est pas un, puisqu'il n'est connu qu'à la clôture. Quand l'espérance d'une \omterm{def:qm:probability-at-speed:martingale}{martingale} en un
\omterm{def:qm:probability-at-speed:stopping}{temps d'arrêt} reste-t-elle égale à sa valeur initiale~?

\begin{definition}[Uniforme intégrabilité]\label{def:qm:probability-at-speed:ui}
Une famille $(X_i)_{i \in I}$ de variables aléatoires possède une \emph{uniforme
intégrabilité}\index{uniforme intégrabilité} si $\sup_{i}\E[|X_i|\mathbf 1_{\{|X_i| > c\}}]
\to 0$ lorsque $c \to \infty$.
\end{definition}

\begin{theorem}[Théorème d'arrêt]\label{thm:qm:probability-at-speed:optional}
Soit $(M_n)$ une \omterm{def:qm:probability-at-speed:martingale}{martingale} et $\tau$ un \omterm{def:qm:probability-at-speed:stopping}{temps d'arrêt}. Alors $\E[M_\tau] = \E[M_0]$ si l'une des deux conditions suivantes est vérifiée~:
(i) $\tau$ est borné, ou (ii) $\tau < \infty$ presque sûrement et le processus arrêté
$(M_{n \wedge \tau})$ est uniformément intégrable~; en particulier si $|M_{n\wedge\tau}| \le c$ pour
tout $n$.
\end{theorem}

\begin{proof}
(i) Si $\tau \le N$, alors $M_\tau = M_0 + \sum_{k=1}^N D_k\mathbf 1_{\{\tau \ge k\}}$, et
$\{\tau \ge k\} = \{\tau \le k-1\}^c \in \mathcal F_{k-1}$, donc chaque terme est de moyenne nulle d'après
la \cref{prop:qm:probability-at-speed:orthogonal}. (ii) Appliquer (i) à $\tau \wedge n$~; $M_{n \wedge
\tau} \to M_\tau$ presque sûrement, et l'\omterm{def:qm:probability-at-speed:ui}{uniforme intégrabilité} améliore cette convergence en
$L^1$ (Vitali), de sorte que les moyennes convergent.
\end{proof}

Le théorème permet de valoriser toute règle de prise de bénéfice et de stop-loss sur une \omterm{def:qm:probability-at-speed:martingale}{martingale}. Une position dont
la valeur varie de $\pm 1$ tick avec une probabilité égale est clôturée à $+a$ ou $-b$. La
valeur arrêtée est bornée, donc $\E[S_\tau] = 0$~: avec $\P(\text{hausse en premier}) = \pi$,
$\pi a - (1 - \pi)b = 0$ et
\[
\pi = \frac{b}{a+b}.
\]
Un stop proche de l'entrée et un objectif de profit éloigné gagnent rarement et perdent souvent, et le
profit espéré est nul quels que soient $a$ et $b$~: la règle remodèle la distribution du
résultat, non sa moyenne. Avec une dérive ($p \ne \tfrac12$), la \omterm{def:qm:probability-at-speed:martingale}{martingale} $(q/p)^{S_n}$ donne
$\pi = (1 - (q/p)^b)/(1 - (q/p)^{a+b})$~; la \cref{fig:qm:probability-at-speed:stops} compare
les deux formules aux résultats de simulation.

\begin{omfigure}
\begin{tikzpicture}
\begin{axis}[width=11cm,height=5cm,
  xlabel={distance de stop-loss $b$ (ticks), prise de bénéfice $a = 10$},ylabel={P(prise de bénéfice en premier)},
  tick label style={font=\small},label style={font=\small},
  xmin=0,xmax=21,ymin=0,ymax=0.75,grid=major,grid style={gray!25},
  xtick={2,4,6,8,10,12,14,16,18,20},
  legend columns=2,legend style={font=\footnotesize,at={(0.5,-0.3)},anchor=north,draw=none,fill=none,/tikz/every even column/.append style={column sep=8pt}},
  table/col sep=comma]
\addplot[omDef,very thick] table[x=b,y=fair_theory]{figdata/methods/01-probability-at-speed/stops.csv};
\addplot[omDef,only marks,mark=*,mark size=1.8pt] table[x=b,y=fair_sim]{figdata/methods/01-probability-at-speed/stops.csv};
\addplot[omThm,very thick] table[x=b,y=drift_theory]{figdata/methods/01-probability-at-speed/stops.csv};
\addplot[omThm,only marks,mark=triangle*,mark size=2.2pt] table[x=b,y=drift_sim]{figdata/methods/01-probability-at-speed/stops.csv};
\legend{{$p = 0.5$: $b/(a+b)$},simulé,{$p = 0.49$: formule},simulé}
\end{axis}
\end{tikzpicture}

\omcaption{Probabilité qu'une marche aléatoire de $\pm1$ tick atteigne une prise de bénéfice à 10 ticks
avant un stop-loss à $b$ ticks~: les formules du théorème d'arrêt (courbes) et 20\,000 trajectoires à graine fixée
par point (marques). Une dérive de 1 point contre la position ($p = 0.49$) coûte 10 points
de probabilité pour $b = 10$. Données~: le tutoriel du chapitre.}\label{fig:qm:probability-at-speed:stops}
\end{omfigure}

\begin{remark}[La stratégie du doublement]\label{rem:qm:probability-at-speed:doubling}
Doubler la mise après chaque perte d'un pari équitable à cote unitaire, et s'arrêter au premier gain
rapporte $+1$ presque sûrement, ce qui semble contredire le théorème. Le \omterm{def:qm:probability-at-speed:stopping}{temps d'arrêt} est fini
mais non borné, et le processus arrêté n'est pas uniformément intégrable~: avant le gain, il a
perdu $2^k - 1$ avec probabilité $2^{-k}$. Une ligne de crédit finie borne les mises, rend
la condition (ii) vérifiée et rétablit $\E[M_\tau] = 0$~: un petit gain fréquent, payé par une
perte rare et énorme.
\end{remark}

\section{Changement de probabilité}\label{sec:qm:probability-at-speed:measure}

Une grande partie de la finance quantitative calcule sous une mesure autre que celle qui décrit le
monde~: une mesure sous laquelle les prix actualisés sont des \omterm{def:qm:probability-at-speed:martingale}{martingales} (chapitre~5), ou une mesure sous laquelle une
perte rare devient assez fréquente pour être simulée (chapitre~26).

\begin{definition}[Changement de probabilité]\label{def:qm:probability-at-speed:change}
Deux mesures de probabilité $\P$ et $\mathbb Q$ sur $(\Omega, \mathcal F)$ sont
des \emph{mesures équivalentes}\index{mesures équivalentes} si elles ont les mêmes ensembles négligeables. Un
\emph{changement de probabilité}\index{changement de probabilité} de $\P$ vers une $\mathbb Q$ équivalente est
décrit par la \emph{dérivée de Radon--Nikodym}\index{dérivée de Radon--Nikodym} $Z =
d\mathbb Q/d\P$, l'unique variable aléatoire positive à un ensemble négligeable près, telle que $\mathbb Q(A) =
\E[Z\mathbf 1_A]$ pour tout $A \in \mathcal F$~; alors $\E^{\mathbb Q}[X] = \E[ZX]$. Le long d'une
\omterm{def:qm:probability-at-speed:filtration}{filtration}, le \emph{processus densité}\index{processus densité} est $Z_t = \E_t[Z]$, la
dérivée de Radon--Nikodym de $\mathbb Q$ restreinte à $\mathcal F_t$.
\end{definition}

Le \omterm{def:qm:probability-at-speed:change}{processus densité} est une $\P$-martingale positive avec $Z_0 = 1$ (une \omterm{def:qm:probability-at-speed:doob}{martingale de Doob}), et
réciproquement, toute \omterm{def:qm:probability-at-speed:martingale}{martingale} de ce type définit une mesure. Les \omterm{def:qm:probability-at-speed:condexp}{espérances conditionnelles} sous la nouvelle
mesure se déduisent des anciennes.

\begin{proposition}[Formule de Bayes pour un changement de probabilité]\label{prop:qm:probability-at-speed:bayes}
Si $\mathbb Q \sim \P$ de \omterm{def:qm:probability-at-speed:change}{processus densité} $(Z_t)$ et si $X$ est $\mathbb Q$-intégrable, alors
pour $s \le t$ et $X$ $\mathcal F_t$-mesurable,
\[
\E^{\mathbb Q}_s[X] = \frac{\E_s[Z_tX]}{Z_s}.
\]
En particulier, un $(X_t)$ adapté est une $\mathbb Q$-martingale si et seulement si $(Z_tX_t)$ est une
$\P$-martingale.
\end{proposition}

\begin{proof}
Le membre de droite $Y$ est $\mathcal F_s$-mesurable. Pour $A \in \mathcal F_s$,
\[
\E^{\mathbb Q}[Y\mathbf 1_A] = \E[Z_sY\mathbf 1_A] = \E\bigl[\E_s[Z_tX]\mathbf 1_A\bigr] =
\E[Z_tX\mathbf 1_A] = \E^{\mathbb Q}[X\mathbf 1_A],
\]
en utilisant $\E[Z W] = \E[Z_s W]$ pour toute variable $\mathcal F_s$-mesurable $W$, et la propriété de la tour.
\end{proof}

\begin{example}[Décaler une gaussienne, et voir une queue lointaine]\label{ex:qm:probability-at-speed:shift}
Soit $X \sim \mathcal N(0, 1)$ sous $\P$ et $Z = \exp(cX - c^2/2)$. Alors $\E^{\mathbb
Q}[e^{\iu uX}] = \E[e^{(c + \iu u)X - c^2/2}] = e^{\iu uc - u^2/2}$~: sous $\mathbb Q$, $X \sim
\mathcal N(c, 1)$. Le \omterm{def:qm:probability-at-speed:change}{changement de probabilité} a déplacé la moyenne sans toucher à la forme. On le parcourt à
rebours pour estimer $p = \P(X > 4) = 3.17 \times 10^{-5}$~: on tire $Y \sim \mathcal N(4, 1)$
et l'on moyenne $e^{-4Y + 8}\mathbf 1_{\{Y > 4\}}$, l'indicatrice repondérée par $d\P/d\mathbb Q$.
Avec 100\,000 tirages, l'échantillonnage naïf ne voit que trois dépassements et a une erreur type de $1.7
\times 10^{-5}$, soit la moitié de la réponse~; l'estimation repondérée a une erreur type de $2.1 \times
10^{-7}$, quatre-vingts fois plus petite. Le chapitre~26 en fait un échantillonnage préférentiel, et le chapitre~5
fait le même calcul pour des trajectoires browniennes entières.
\end{example}

\section{Notations de la collection}\label{sec:qm:probability-at-speed:notation}

Chaque livre de la collection, à partir de celui-ci, utilise les symboles ci-dessous. Ils prolongent les tableaux des
livres~1 et~2, dont les sens sont conservés~: $s_t$ est la fourchette, $\mathcal S$
un spread de crédit, $\tau = T - t$ le temps restant jusqu'à l'échéance, $r$ un taux, $\ell$ des frais d'emprunt, $K$ un
prix d'exercice, $P(t,T)$ un facteur d'actualisation, $y$ un rendement, $\delta$ une fraction de couru et $R$ un
taux de recouvrement. Un symbole qu'un chapitre déclare \emph{local} peut y avoir un autre sens.

\begin{notation}[Probabilités, mesures et processus]\label{not:qm:probability-at-speed:prob}
\leavevmode\par\smallskip\noindent
{\small
\begin{tabular}{@{}p{4.3cm}p{9.7cm}@{}}
\toprule
$(\Omega,\mathcal F,\P)$, $\mathbb F = (\mathcal F_t)$ & espace de probabilité~; \omterm{def:qm:probability-at-speed:filtration}{filtration} (conditions usuelles) \\
$\E_t[X] = \E[X \mid \mathcal F_t]$, $\E^{\mathbb Q}_t$ & \omterm{def:qm:probability-at-speed:condexp}{espérance conditionnelle}~; sous une autre probabilité \\
$\P$, $\mathbb Q$, $B_t = e^{\int_0^t r_s\,ds}$ & probabilité historique~; probabilité risque-neutre~; compte monétaire \\
$\mathcal N_t$, $\mathbb Q^{\mathcal N}$~; $\mathbb Q^T$, $\mathbb Q^A$ & numéraire générique et sa probabilité~; probabilités $T$-forward et annuité \\
$d\mathbb Q/d\P$, $Z_t$ & \omterm{def:qm:probability-at-speed:change}{dérivée de Radon--Nikodym}~; \omterm{def:qm:probability-at-speed:change}{processus densité} \\
$\Phi$, $\varphi$~; $\mathcal N(m, s^2)$ & fonction de répartition et densité de la loi normale centrée réduite (jamais $N(d_1)$)~; la loi normale, toujours avec ses arguments \\
$\varphi_X(u) = \E[e^{\iu uX}]$ & \omterm{def:qm:probability-at-speed:charfn}{fonction caractéristique}, toujours indexée \\
$\mathbf 1_A$~; $\overset{d}{=}$, $\xrightarrow{d}$, $\xrightarrow{\P}$ & indicatrice~; égalité et convergence en loi, en probabilité \\
$\tau$ & un \omterm{def:qm:probability-at-speed:stopping}{temps d'arrêt}~; là où il désigne un temps restant jusqu'à l'échéance, écrire $T - t$~; les instants de défaut sont indexés ($\tau_C$) \\
$W_t$~; $W^{\mathbb Q}_t$ & mouvement brownien sous la probabilité en vigueur~; muni d'un exposant lorsque deux probabilités apparaissent \\
$[X]_t$, $[X,Y]_t$, $\mathcal E(X)_t$ & variation quadratique, covariation, exponentielle stochastique \\
$dX = \mu\,dt + \sigma\,dW$~; $\mathcal L$, $\mathcal L^*$ & une diffusion~; son générateur et son adjoint \\
$dX = \kappa(\bar x - X)\,dt + \sigma\,dW$ & processus d'Ornstein--Uhlenbeck~: vitesse $\kappa$, niveau $\bar x$, demi-vie $\ln 2/\kappa$ \\
$dv = \kappa(\bar v - v)\,dt + \eta\sqrt v\,dW$ & processus en racine carrée~; condition de Feller $2\kappa\bar v \ge \eta^2$~; $\eta$ désigne partout la volatilité de la volatilité \\
$N_t$, $t_1 < t_2 < \dots$ & processus de comptage et ses instants d'événement \\
$\lambda_t$, $\Lambda_t = \int_0^t\lambda_s\,ds$ & intensité (taux de Poisson, intensité de défaut, intensité de Hawkes $\lambda_t = \mu + \sum_{t_i < t}g(t - t_i)$)~; compensateur \\
$(\sigma^2, \nu, \gamma)$, $\psi$ & triplet de Lévy et exposant caractéristique, $\varphi_{X_t}(u) = e^{t\psi(u)}$ \\
$H$, $W^H_t$ & exposant de Hurst~; mouvement brownien fractionnaire \\
\bottomrule
\end{tabular}}
\end{notation}

\begin{notation}[Statistique, matrices et calcul numérique]\label{not:qm:probability-at-speed:stat}
\leavevmode\par\smallskip\noindent
{\small
\begin{tabular}{@{}p{4.3cm}p{9.7cm}@{}}
\toprule
$n$, $\theta \in \Theta$, $\theta_0$ & taille d'échantillon, paramètre, vraie valeur~; chapeau pour un estimateur, tilde pour un estimateur rétréci, barre pour une moyenne empirique \\
$\ell_n(\theta)$, $\mathcal I(\theta)$ & log-vraisemblance (toujours avec indice et argument)~; information de Fisher \\
$\mathrm{se}(\hat\theta)$~; $\mathrm{SR}$, $\widehat{\mathrm{SR}}$ & erreur type~; ratio de Sharpe et son estimation \\
$R_t = \ln(S_t/S_{t-1})$ & rendement logarithmique ($R$ seul reste le taux de recouvrement) \\
$L$, $\Delta = 1 - L$~; $\gamma(h)$, $\rho(h)$ & opérateurs de retard et de différence~; autocovariance, autocorrélation \\
$\phi_i$, $\vartheta_j$, $\varepsilon_t$ & coefficients autorégressifs et de moyenne mobile (toujours indexés)~; innovations \\
$\mathrm{RV}_t$, $\mathrm{IV}_t$~; $\sigma_{\mathrm{imp}}$ & variance réalisée et variance intégrée~; volatilité implicite (jamais «~IV~» dans une formule) \\
$\Sigma$, $C$, $\hat\Sigma$~; $\mathbf 1$, $I_n$, ${}^\top$ & covariance, corrélation, covariance empirique~; vecteur des uns, matrice identité, transposée~; les vecteurs sont des colonnes \\
$\lambda_1 \ge \dots \ge \lambda_N$, $\operatorname{cond}(A)$ & valeurs propres (toujours indexées)~; conditionnement \\
$\min f(x)$ s.c.\ $g_i(x) \le 0$, $h_j(x) = 0$ & optimisation~; multiplicateurs $u \ge 0$, $v$~; optimum $x^\star$, valeurs primale $p^\star$ et duale $d^\star$ \\
$\mathrm{fl}(x)$, $u = 2^{-53}$, $\varepsilon_{\mathrm{mach}} = 2^{-52}$ & virgule flottante en binaire64, arrondi au plus proche~; $\mathrm{ulp}(x)$ \\
$t_k = k\Delta t$, $x_j$, $V^k_j$~; $M$, $\hat V_M$ & grille temporelle, grille spatiale, solution numérique~; trajectoires de Monte-Carlo et estimateur \\
\bottomrule
\end{tabular}}
\end{notation}

Les symboles des options, des taux, du crédit et du risque ($\sigma_{\mathrm{imp}}(K,T)$, $k = \ln(K/F_{0,T})$,
$\xi_t(u)$, les grecques, la grecque de taux étant notée Rho car $\rho$ désigne toujours une corrélation,
$S_{a,b}(t)$, $A_{a,b}(t)$, $Q_C(t)$, $\mathrm{VaR}_{\alpha,h}$, $\mathrm{ES}_{\alpha,h}$) viennent
de ce même tableau et sont introduits dans les livres~5 et~6.

\section{Tutoriel~: la martingale de Doob d'un jeu de cartes}\label{tut:qm:probability-at-speed:cards}

\begin{tutorial}
\textbf{Objectif.} Construire la \omterm{def:qm:probability-at-speed:doob}{martingale de Doob} du jeu à cinq cartes du livre~2,
chapitre~30 (la somme de cinq cartes tirées d'un jeu de 52 cartes, les as valant 1 et les rois 13), vérifier que
ses révisions se comportent comme le dit la \cref{prop:qm:probability-at-speed:orthogonal}, et repérer un
prévisionniste qui réagit trop peu. \textbf{État final~:} le tableau ci-dessous et
les \cref{fig:qm:probability-at-speed:day,fig:qm:probability-at-speed:stops}.
\begin{tutsteps}
\item \textbf{La prévision.} Après $k$ cartes de somme $s_k$, les $52 - k$ cartes restantes
ont pour somme $364 - s_k$, donc $M_k = s_k + (5 - k)(364 - s_k)/(52 - k)$, et $M_0 = 35$. Le
prévisionniste qui réagit trop peu ne parcourt qu'une fraction $a$ du chemin depuis 35 jusqu'à la dernière carte.
\omcode{code/methods/01-probability-at-speed/python/qm_martingales.py}{17}{39}{La martingale de Doob du jeu de cartes et un prévisionniste qui réagit trop peu.}\label{lst:qm:probability-at-speed:cards}
\item \textbf{Le test.} Régresser la dernière révision sur la précédente, en regroupant les donnes~;
la fonction du projet en cours le fait pour tout panel de trajectoires de prévision.
\omcode{code/firm/mgtest/firm_mgtest.py}{29}{47}{Régression d'une révision sur une révision antérieure~: pente nulle pour une martingale.}\label{lst:qm:probability-at-speed:regression}
\item \textbf{Exécuter} \texttt{card\_table()} sur 20\,000 donnes à graine fixée, puis
\texttt{fig\_martingales.py} pour les figures.
\end{tutsteps}
\begin{center}
\small
\begin{tabular}{@{}lcc@{}}
\toprule
 & prévision honnête & sous-réaction ($a = 0.6$) \\
\midrule
révision moyenne & $0.001$ & $0.001$ \\
pente dernière révision / quatrième & $-0.005$ & $0.661$ \ (théorie $0.667$) \\
statistique $t$ de la pente & --- & $45.9$ \\
\bottomrule
\end{tabular}
\end{center}
Les deux prévisions ont des révisions de moyenne nulle~; seule la régression les distingue.
\textbf{Pour aller plus loin.} Faire surréagir le prévisionniste ($a > 1$) et prévoir le signe de
la pente~; remplacer le jeu de cartes par le modèle du cours de clôture du problème du week-end et déterminer
combien de jours de données le test requiert.
\end{tutorial}


\section{Construction~: diagnostics de martingale}\label{bld:qm:probability-at-speed:mgtest}

\begin{build}
\bfield{Rôle} Toute prévision que publie la petite société (une fair value, une espérance du
cours de clôture, un taux d'exécution prévu) est vérifiée au regard de la propriété de \omterm{def:qm:probability-at-speed:martingale}{martingale} avant que personne
n'opère sur elle ou contre elle.
\bfield{Interface} \texttt{revisions(paths)}~; \texttt{revision\_regression(paths, lag=1,
col=None)} renvoyant la pente, l'erreur type, $t$ et $n$~; \texttt{variance\_ratio(increments,
q)}~; \texttt{variance\_shares(paths)}~; \texttt{martingale\_report(paths)}. Un panel a une
ligne par jour ou par donne et une colonne par date de révision.
\bfield{Règles} Pas d'ordonnée à l'origine dans la régression des révisions (la révision moyenne est rapportée
séparément)~; les parts de variance sont celles des révisions, qui, pour une \omterm{def:qm:probability-at-speed:martingale}{martingale}, somment à la
variance de la valeur finale~; tout est à graine fixée et vectorisé.
\bfield{Tests d'acceptation} \texttt{code/firm/mgtest/tests/}~: les marches aléatoires passent~; une prévision
qui intègre tardivement la moitié des nouvelles est repérée avec la bonne pente~; les parts sont uniformes pour des
révisions iid~; le ratio de variance vaut un pour des accroissements iid et est inférieur à un pour des
accroissements à retour à la moyenne.
\bfield{Pour aller plus loin} Erreurs types robustes à l'hétéroscédasticité et à l'autocorrélation (chapitre~11)~;
un test des révisions contre tout signal $\mathcal F_{k-1}$-mesurable, et non seulement des
révisions passées~; contrôle des tests multiples sur de nombreuses prévisions (chapitre~12).
\end{build}

\begin{omsources}
\item A.~N.~Kolmogorov, \emph{Grundbegriffe der Wahrscheinlichkeitsrechnung}, Springer,
1933~: les axiomes de la théorie de la mesure.
\item J.~Ville, \emph{Étude critique de la notion de collectif}, Gauthier-Villars, 1939,
où le mot martingale fait son entrée en probabilités.
\item J.~L.~Doob, \emph{Stochastic Processes}, Wiley, 1953.
\item O.~Nikodym, ``Sur une généralisation des intégrales de M.~J.~Radon'', \emph{Fundamenta
Mathematicae} 15, 1930.
\item D.~Williams, \emph{Probability with Martingales}, Cambridge University Press, 1991~: les
démonstrations admises ici.
\end{omsources}

\section{Exercices}

\begin{exercise}[$\star$]\label{exo:qm:probability-at-speed:1}
Dans le jeu à cinq cartes, calculer la prévision de la somme une fois que la première carte
s'avère être un roi.
\end{exercise}

\begin{exercise}[$\star$]\label{exo:qm:probability-at-speed:2}
Indiquer lesquelles de ces dates sont des \omterm{def:qm:probability-at-speed:stopping}{temps d'arrêt} pour la \omterm{def:qm:probability-at-speed:filtration}{filtration} des cours observés~: (a) la première date
où le cours dépasse de 1~\% celui de l'ouverture~; (b) l'instant du plus bas de la journée~; (c) 15~h~50~; (d) la première
date postérieure à 15~h~50 où la prévision de la clôture varie de plus de 10 centimes~?
\end{exercise}

\begin{exercise}[$\star$]\label{exo:qm:probability-at-speed:3}
Une position sur une marche équitable à $\pm1$~tick est clôturée en $+1$ ou en $-3$. Calculer la
probabilité de la prise de profit, et expliquer pourquoi le profit espéré reste nul.
\end{exercise}

\begin{exercise}[$\star\star$]\label{exo:qm:probability-at-speed:4}
Une prévision de la clôture est révisée par des nouvelles indépendantes, de même variance, à chacune des 390
minutes d'une séance, et rien d'autre. Calculer la part de la variance de la clôture qui est résolue
dans la dernière demi-heure.
\end{exercise}

\begin{exercise}[$\star\star$]\label{exo:qm:probability-at-speed:5}
Sous $\P$, $X \sim \mathcal N(0,1)$, et $d\mathbb Q/d\P = \exp(cX - c^2/2)$. Calculer $\E^{\mathbb
Q}[X]$ et $\mathbb Q(X > c)$, et vérifier que $\mathbb Q$ est une mesure de probabilité.
\end{exercise}

\begin{exercise}[$\star\star$]\label{exo:qm:probability-at-speed:6}
Le prévisionniste sous-réactif du jeu de cartes publie $F_k = 35 + a(M_k - 35)$ pour $k \le 4$ et $F_5 =
M_5$. Calculer $\E[F_5 - F_4 \mid \mathcal F_4]$ et la pente de la régression de $F_5 - F_4$ sur
$F_4 - 35$.
\end{exercise}

\begin{exercise}[$\star\star\star$]\label{exo:qm:probability-at-speed:7}
\emph{Programmation.} Avec \texttt{revision\_regression}, estimer la pente de la dernière révision
du prévisionniste sous-réactif du jeu de cartes ($a = 0.6$) sur la révision précédente, sur 20\,000
donnes à graine fixée, et la comparer à la réponse de l'\cref{exo:qm:probability-at-speed:6}.
\end{exercise}

\begin{exercise}[$\star\star\star$]\label{exo:qm:probability-at-speed:8}
\emph{Repérer l'erreur.} «~Sur un an, les révisions de notre modèle de fair value ont été nulles en moyenne à quatre
décimales près, donc le modèle est une \omterm{def:qm:probability-at-speed:martingale}{martingale} et personne ne peut prendre position contre lui.~» Corriger ce raisonnement.
\end{exercise}

\section{Problème~: la fair value de la clôture}

\begin{problem}[{Problème du week-end --- quelle part de la clôture est encore inconnue à 15~h~50}]\label{pb:qm:probability-at-speed:1}
Un desk publie $M_t = \E_t[C]$, sa prévision du cours de clôture $C$ d'une action. Dans son modèle,
le cours part de 100~USD et évolue à chacune des 390 minutes de la séance par
accroissements gaussiens indépendants, d'écart type total de 1.80~USD sur la journée~; à
15~h~50, le déséquilibre de l'enchère de clôture est publié~; il déplace la prévision
d'un montant gaussien indépendant $J$, d'écart type 0.30~USD et de moyenne nulle~; les dix
dernières minutes de la séance suivent.

\medskip\noindent\textbf{Partie I --- La prévision.}
\begin{enumerate}
\item Calculer la variance et l'écart type de $C$.
\item Écrire $M_t$ avant 15~h~50 et après l'annonce, en fonction du cours observé.
\item Expliquer pourquoi $(M_t)$ est une \omterm{def:qm:probability-at-speed:martingale}{martingale} pour la \omterm{def:qm:probability-at-speed:filtration}{filtration} du desk.
\item Calculer la révision de $M$ à l'annonce de 15~h~50.
\item Calculer la probabilité que l'annonce déplace la prévision de plus de 14 centimes.
\end{enumerate}

\medskip\noindent\textbf{Partie II --- Ce qui est résolu et quand.}
\begin{enumerate}[resume]
\item Expliquer pourquoi les variances des révisions s'additionnent pour donner $\Var(C)$.
\item Calculer la part de $\Var(C)$ résolue à midi.
\item Calculer la part résolue à l'annonce elle-même.
\item Calculer la part résolue de l'annonce à la clôture, bornes incluses.
\item Vérifier cette part par rapport à la simulation de 20\,000 journées.
\end{enumerate}

\medskip\noindent\textbf{Partie III --- Une prévision trop lente.}
\begin{enumerate}[resume]
\item Un second desk n'intègre que la moitié de $J$ à 15~h~50 et l'autre moitié à la clôture. Dire si
sa prévision est une \omterm{def:qm:probability-at-speed:martingale}{martingale}.
\item Calculer la corrélation entre sa révision de 15~h~50 et sa révision sur les dix
dernières minutes.
\item Calculer la pente que donne la régression de la seconde révision sur la première, en théorie et
dans la simulation.
\item Déterminer le nombre de journées de données nécessaire pour que la statistique $t$ de la pente atteigne 2.
\item Un trader achète une action à ce desk au prix de sa prévision à 15~h~50 lorsque sa révision est
positive, en vend une lorsqu'elle est négative, et clôture sa position à la clôture. Calculer le
gain moyen par action.
\end{enumerate}

\medskip\noindent\textbf{Partie IV --- Jugement.}
\begin{enumerate}[resume]
\item Dans l'hypothèse où les dix dernières minutes portent une variance par minute triple de celle du reste de la
journée (le total de la journée restant inchangé), calculer la part résolue de l'annonce à la clôture.
\item Dire si la réponse dépend du niveau du cours.
\item Expliquer pourquoi cette part compte pour un trader qui ne peut agir que dans l'enchère de clôture.
\item Énoncer le \emph{résultat nommé}~: la part de la variance de la clôture résolue à partir de 15~h~50.
\item En une phrase, dire ce que garantit la propriété de \omterm{def:qm:probability-at-speed:martingale}{martingale} d'une prévision, et ce
qu'elle ne garantit pas.
\end{enumerate}
\end{problem}

\section{Questions d'entretien}

\begin{interviewq}[$\star$ \iqroles{trader}]\label{iq:qm:probability-at-speed:1}
Vous jouez à un jeu de pile ou face équilibré, à un dollar la mise, et vous vous arrêtez dès que vous êtes en avance d'un dollar.
N'est-ce pas un gain garanti~?
\end{interviewq}

\begin{interviewq}[$\star$ \iqroles{chercheur}]\label{iq:qm:probability-at-speed:2}
Qu'est-ce qu'une \omterm{def:qm:probability-at-speed:martingale}{martingale}~? Un cours d'action en est-il une~?
\end{interviewq}

\begin{interviewq}[$\star\star$ \iqroles{chercheur, ingénieur ML}]\label{iq:qm:probability-at-speed:3}
Montrer que les accroissements d'une \omterm{def:qm:probability-at-speed:martingale}{martingale} sont non corrélés. Qu'est-ce que cela
implique pour une prévision que vous publiez~?
\end{interviewq}

\begin{interviewq}[$\star\star$ \iqroles{trader}]\label{iq:qm:probability-at-speed:4}
Au moment de la publication du déséquilibre, votre prévision du cours de clôture bouge de 14 cents~; celle d'un concurrent
bouge de 7 cents, puis de 7 autres à la clôture, chaque jour. Que faites-vous~?
\end{interviewq}

\begin{interviewq}[$\star\star$ \iqroles{chercheur, risque}]\label{iq:qm:probability-at-speed:5}
Quand $\E[X \mid Y]$ est-elle la régression linéaire de $X$ sur $Y$~?
\end{interviewq}

\begin{interviewq}[$\star\star\star$ \iqroles{développeur, chercheur}]\label{iq:qm:probability-at-speed:6}
Concevoir un contrôle quotidien vérifiant qu'un service de juste valeur en temps réel se comporte comme une \omterm{def:qm:probability-at-speed:martingale}{martingale}. Qu'est-ce qui peut
rendre ce contrôle trompeur~?
\end{interviewq}
